\documentclass[11pt, oneside]{article}   	% use "amsart" instead of "article" for AMSLaTeX format
\usepackage[margin = 1in]{geometry}                		% See geometry.pdf to learn the layout options. There are lots.
\geometry{letterpaper}                   		% ... or a4paper or a5paper or ... 
%\geometry{landscape}                		% Activate for rotated page geometry
\usepackage[parfill]{parskip}    		% Activate to begin paragraphs with an empty line rather than an indent
\usepackage{graphicx, array,hyperref}				% Use pdf, png, jpg, or eps§ with pdflatex; use eps in DVI mode
								% TeX will automatically convert eps --> pdf in pdflatex		
\usepackage{amssymb, enumerate, tikz, multicol}


\newcommand{\be}{\begin{enumerate}}
\newcommand{\ee}{\end{enumerate}}


\title{Math F113X: Homework Set 11}
%\author{The Author}
\date{}							% Activate to display a given date or no date

\begin{document}

\maketitle

\vspace{-1.5cm}


\fbox{\parbox{\textwidth}{
All of the problems come from the Math in Society text found here:\\
\url{https://spot.pcc.edu/math/mathinsociety/root-1-2.html}\\

\begin{itemize}
\item Complete problems \#11,12,13  from Section 2.1.6 linked here:\\
\url{https://spot.pcc.edu/math/mathinsociety/chap_two_intro_spreadsheet.html}\\

\item Complete problems \#1-6,8,10 from Section 2.2.11 linked here:\\
\url{https://spot.pcc.edu/math/mathinsociety/chap_two_simple_compound.html}\\


\item Complete problems \#1,2,3,4,5 from Section 2.4.9 linked here:\\
\url{https://spot.pcc.edu/math/mathinsociety/chap_two_loans.html}\\


\item Problem A (below)
\end{itemize}
}}

\fbox{Problem A} You have a credit card bill of \$4923. You decide that's way too much money to have floating on your credit card, which has a terrible APR of 27.99\%. So, you \emph{stop using the credit card}.

\begin{enumerate}[(a)]
\item How much interest is charged in the next month, if you make no payments?
\item Why is paying off your credit card at \$100/month not a good strategy?
\item If you pay \$400/month, how many years will it take to pay down the credit card?
\item If you want to pay off the credit card completely in 2 years, how much do you have to pay each month?
\end{enumerate}

\hrulefill

Remember to write up your homework solutions according to the homework writeup guidelines. 

Homework is graded using the following rubric for each problem (or problem part):

\begin{description}
\item[2 points:] You provided a complete answer, with supporting work, written up clearly
\item[1 point:] Some attempt at a solution, but incomplete writeup / unclear / illegible / no answer
\item[1 point:] Only an answer, with no supporting work 
\item[0 points:] Missing.
\end{description}

After you do the homework, you need to check your answers against the solutions! Then figure out your errors (if any) and revise your homework before you submit it. 

Homework must be submitted on Gradescope.

\end{document}  